\documentclass[11pt,letterpaper]{article}
\usepackage[utf8]{inputenc}
\usepackage[spanish]{babel}
\usepackage{amsmath}
\usepackage{amsfonts}
\usepackage{amssymb}
\usepackage[left=1in,right=1in,top=1in,bottom=1in]{geometry}
%Paquete para graficar
\usepackage{pgfplots}
\usepackage{multicol}
\usepackage{color}
\pgfplotsset{compat=1.8}
\usepackage{enumerate}

%Tiks
\usepackage{physics}
\usepackage{amsmath}
\usepackage{tikz}
\usepackage{mathdots}
\usepackage{yhmath}
\usepackage{cancel}
\usepackage{color}
\usepackage{siunitx}
\usepackage{array}
\usepackage{multirow}
\usepackage{amssymb}
\usepackage{gensymb}
\usepackage{tabularx}
\usepackage{extarrows}
\usepackage{booktabs}
\usetikzlibrary{fadings}
\usetikzlibrary{patterns}
\usetikzlibrary{shadows.blur}
\usetikzlibrary{shapes}



% PAQUETES %%%%%%%%%%%%%%%%%%%%%%%%%%%%%%%%%%%%%%%%%%%%%%%%%%%%%%%%%%%%%%%%%%%%

\usepackage[utf8]{inputenc}
\usepackage{graphicx}
\usepackage{listings}
\usepackage{xcolor}
\usepackage{multicol}
\usepackage{anysize} %Permitir distintas medidas de margenes
\usepackage{framed}
\usepackage{fancyhdr}
\usepackage{float}


\fancyhf{}
\rhead{ING2104 - Fundamentos de Economía\\2022-10\vspace{0.1cm}}
\lhead{\includegraphics[scale=0.18]{ing.png}}
\lfoot{\scriptsize abanduc@miuandes.cl}
\rfoot{\thepage}

\renewcommand{\footrulewidth}{0.25pt}

\addtolength{\headheight}{1.4cm}


\newcommand{\evaluationtitle}[2]{
	% T\'itulo
	\begin{center}
	\vspace{1ex}\Large #1\\
	\vspace{1ex}\small #2
	\end{center}
}


\pagestyle{fancy}

\begin{document}


\begin{center}
    \textbf{\Large{Ayudantía 12}}\\
    Ayudantes: Antonia Banduc, Vicente Muñoz, Hernan Venegas, Roberto Witting
\end{center}


\section*{Ejercicio 1}
Imagine que un amigo le pide prestado \$80.000 y le devuelve \$95.000 a final de año. Suponga que la tasa de inflación ese año fue de 10\%.
\begin{enumerate}
    \item Calcule la tasa de interés nominal que su amigo le pagaría.
\item Calcule la tasa de interés real que su amigo le pagaría.
\item Si su amigo se demora dos años en pagarle y mantiene la tasa de interés, ¿Cuánto recibiría usted?
\item  Usted le presta el dinero a su amigo en UF, cuando vale \$30.000, si su amigo mantiene la tasa de
interés real y el prestamo es por un año, ¿Cuánto dinero recibe?
\item Ahora suponga que su amigo le ofrece devolverle \$110.000 al finalizar el tercer año. ¿Conviene darle
el préstamo?
\end{enumerate}

\section*{Ejercicio 2}
Suponga que hace un año hizo un depósito de \$500.000 en una cuenta, justo un año después usted tiene
\$750.000.
\begin{enumerate}
    \item ¿Cuál fue la tasa de interés nominal mensual que obtuvo?
    \item ¿Cuál fue la tasa de interés real que obtuvo, si la inflación fue de 5\%?
\end{enumerate}

\section*{Ejercicio 3}
Si hago un depósito de \$100.000 en una cuenta que entrega un interés de 10\% anual, ¿Cuánto dinero habrá en la cuenta después de 3 años?

\section*{Ejercicio 4}
Si tengo \$250.000 y quiero que se transformen en \$320.000 en un fondo que entrega 2.5\% de rentabilidad anual, ¿Cuánto tiempo debo esperar?


\end{document}